%=====================================================================
\section{Main results of version 0.1.5}\label{sec:main}

Throughout, $M(t)=\|\omega(t)\|_{L^\infty}$ is \lean{vorticity_sup_norm (vorticity (u t))}, and
\lean{IsPeriodicSmoothSolutionOn ν T u₀ u p} is the class of classical solutions on $[0,T)$: $(u,p)$ jointly
$C^\infty$ on $[0,T)\times\R^3$, the equations with one-sided time derivative and zero forcing,
$\operatorname{div}u=0$, $u(0)=u_0$, and $u(t),p(t)$ $\mathbb{Z}^3$-periodic (exact text in
Appendix~\ref{app:lean:clay}).

\begin{definition}[The Frohmanian tether bound; \lean{TetherImpliesBKM.PeriodicTetherBound}]\label{def:ptb}
For $\nu\in\R$ and a field $u_0$, \lean{PeriodicTetherBound ν u₀} asserts that there is $Y\in\R$ such that
$M(t)\le Y$ for every $T$, every classical periodic solution $(u,p)$ from $u_0$ on $[0,T)$ and every
$t\in[0,T)$. The \emph{closing bound} \lean{PeriodicTetherClosingBound ν} asserts \lean{PeriodicTetherBound ν u₀}
for every smooth divergence-free $\mathbb{Z}^3$-periodic $u_0$; this is a definitional identity
(\lean{periodicTetherClosingBound_iff}, proved by \lean{Iff.rfl}).
\end{definition}

\begin{definition}[Periodic Beale--Kato--Majda integrability; \lean{TetherImpliesBKM.PeriodicBKMIntegrability}]
\lean{PeriodicBKMIntegrability ν u₀} asserts: for every horizon $T_0\in\R$ there is $C\in[0,\infty)$ such
that
\[
\int_{[0,T)} M(t)\,dt\;\le\;C
\]
for every $T\le T_0$ and every classical periodic solution from $u_0$ on $[0,T)$. The integral is the lower
Lebesgue integral $\int^{-}_{[0,T)}\mathrm{ofReal}(M(t))\,dt\in[0,\infty]$ (Mathlib's \lean{lintegral} with
\lean{ENNReal.ofReal}), so no measurability side condition is involved.
\end{definition}

\begin{theorem}[Main theorem, kernel-checked; \lean{ClayStatements.clay_B_exact_of_periodic_tether_bound}]\label{thm:mainA}
Assume, for every $\nu>0$,
\begin{enumerate}[label=(\roman*)]
\item \lean{PeriodicTetherClosingBound ν} (the Frohmanian tether bound, Definition~\ref{def:ptb}), and
\item \lean{PeriodicBKMContinuation ν}: for every smooth divergence-free $\mathbb{Z}^3$-periodic $u_0$, a single
ceiling $Y$ for $M$ along all classical periodic solutions from $u_0$ on all intervals $[0,T)$ yields a global
solution $(u,p)$ satisfying \lean{IsPeriodicSmoothSolution ν u₀ (fun _ _ => 0) u p}.
\end{enumerate}
Then \lean{Clay_B_exact} holds: for every $\nu>0$ and every smooth, divergence-free, $\mathbb{Z}^3$-periodic
$u_0$ there exist $(u,p)$, jointly $C^\infty$ on $[0,\infty)\times\R^3$, solving the Navier--Stokes equations
with zero forcing, with $u(0)=u_0$ and $u(t),p(t)$ periodic for all $t\ge0$.
\end{theorem}

Hypothesis (ii) packages three classical results that are not formalized: local existence of classical
periodic solutions, their uniqueness, and the continuation of a solution whose vorticity stays bounded
(Beale--Kato--Majda \cite{BKM1984}). Its antecedent is literally the body of (i); the proof of
Theorem~\ref{thm:mainA} is therefore the composition of (ii) with (i), and the analytic content lies entirely in
the two hypotheses. Hypothesis (i) is the key estimate. It is open; it is not supplied by the tether kernel of
Sections~\ref{sec:tether}--\ref{sec:damping}; and it is as strong as statement (B) itself (for $\nu>0$, global
smoothness on $\T$ together with the decay of the energy gives a time-uniform bound on $M$).

\begin{theorem}[Tether bound $\Rightarrow$ BKM integrability, kernel-checked; \lean{TetherImpliesBKM.periodicBKMIntegrability_of_periodicTetherBound}]\label{thm:mainB}
For every $\nu\in\R$ and every field $u_0$, \lean{PeriodicTetherBound ν u₀} implies
\lean{PeriodicBKMIntegrability ν u₀}. Explicitly, if $M\le Y$ along every classical periodic solution from
$u_0$, then for every $T_0$ the constant $C=\max(Y,0)\cdot\max(T_0,0)$ bounds $\int_{[0,T)}M(t)\,dt$ for every
such solution on $[0,T)$ with $T\le T_0$.
\end{theorem}

\begin{proof}
For $t\in[0,T)$, $\mathrm{ofReal}(M(t))\le\mathrm{ofReal}(\max(Y,0))$, so by monotonicity of the lower
integral on the measurable set $[0,T)$ (\lean{setLIntegral_mono'}), \lean{setLIntegral_const} and
\lean{Real.volume_Ico}, the integral is at most $\mathrm{ofReal}(\max(Y,0))\cdot\mathrm{ofReal}(T)\le
\mathrm{ofReal}(\max(Y,0)\max(T_0,0))<\infty$.
\end{proof}

\begin{corollary}[\lean{TetherImpliesBKM.periodicBKMIntegrability_of_periodicTetherClosingBound}]
\lean{PeriodicTetherClosingBound ν} implies \lean{PeriodicBKMIntegrability ν u₀} for every smooth
divergence-free $\mathbb{Z}^3$-periodic $u_0$.
\end{corollary}

\subsection{The converse direction: status (not a theorem of this paper)}\label{sec:converse}
For $\nu>0$ and a fixed smooth divergence-free periodic $u_0$, the converse implication
\lean{PeriodicBKMIntegrability ν u₀} $\Rightarrow$ \lean{PeriodicTetherBound ν u₀} is true classical mathematics,
by the following route; it is \emph{not} proved in Lean, and no statement of this paper asserts it.
\begin{enumerate}[label=(P\arabic*)]
\item \emph{Beale--Kato--Majda continuation} \cite{BKM1984}: if $\int_0^{T}M\,dt<\infty$ along the maximal
classical solution, the solution extends beyond $T$; with BKM integrability on every horizon the solution is
global.
\item \emph{Uniqueness of classical periodic solutions} (energy estimate for the difference and Gronwall), so
that every classical solution on $[0,T)$ is the restriction of the global one.
\item \emph{Long-time decay}: by the energy inequality $\int_0^\infty\|\nabla u\|_{L^2}^2\,dt<\infty$, so
$\|\nabla u(t_0)\|_{L^2}$ is small at some $t_0$; small-data global theory and parabolic smoothing then give
$\sup_{t\ge t_0+1}M(t)<\infty$, and continuity bounds $M$ on $[0,t_0+1]$. (A nonzero mean velocity is removed by
a Galilean change of frame, which does not change $\omega$.)
\end{enumerate}
These three pillars are the formalization targets of the follow-up work; (P2) is the first, and it also discharges
the uniqueness input contained in \lean{PeriodicBKMContinuation}. With (P1)--(P3) the two bounds are equivalent
datum by datum. Over all data the two statements are each equivalent to statement (B), so the equivalence
locates the tether bound exactly at Beale--Kato--Majda strength; it does not reduce the open problem.

\paragraph{Boundedness of $\omega(t)$ in the periodic class.} For $t\in[0,T)$ a classical periodic solution has
$\omega(t)$ continuous and periodic, hence bounded, so the real supremum \lean{vorticity_sup_norm} takes its genuine
value there. This classical fact is not formalized and is used by no theorem.

%=====================================================================
\section{Corrections in version 0.1.5: the local tether hypotheses}\label{sec:corr}

The definitions of the $\R^3$-class local hypotheses in version 0.1.4 had two defects, found by a definition
audit and now proved in Lean (module \lean{TetherDefinitionAudit}, which keeps the 0.1.4 definitions verbatim
as \lean{LocalTetherBoundV014} and \lean{LocalTetherRateV014}). Both are corrected; the exact statement changes
are listed in Section~\ref{sec:diff}.

\subsection{The junk value of the supremum}
\lean{vorticity_sup_norm ω} is the real supremum $\bigsqcup_x\|\omega(x)\|$, which Mathlib defines to be $0$
when the range is not bounded above.

\begin{theorem}[\lean{vorticity_sup_norm_eq_zero_of_not_bddAbove}, \lean{localTetherBoundV014_constraint_free_of_unbounded}]
If $x\mapsto\|\omega(x)\|$ is unbounded then \lean{vorticity_sup_norm ω} $=0$. Consequently, at every time $t$
at which $\omega(t)$ is unbounded, the 0.1.4 ceiling $M(t)\le Y$ holds for every $Y\ge0$.
\end{theorem}

So the 0.1.4 \lean{LocalTetherBound} and \lean{LocalFiniteTimeVorticityBound} imposed no constraint at such
times. Smooth divergence-free finite-energy data with unbounded vorticity exist: with a bump $\psi\in C_c^\infty$
equal to $1$ near $0$ and centres $c_n$ at mutual distance larger than twice the diameter of its support, let
$A=\sum_n 8^{-n}\sin(4^nx_1)\,\psi(x-c_n)\,e_3$ and $u_0=\nabla\times A$. The sum is locally finite, so $u_0$ is
smooth and divergence-free; $\|u_0\|_{L^2}^2\lesssim\sum_n(8^{-n}4^n)^2<\infty$; and near $c_n$,
$\omega_0=\nabla\times\nabla\times A=-\Delta A=8^{-n}16^n\sin(4^nx_1)\,e_3$ reaches size $2^n$. (This
construction is classical and not formalized.) Every such datum admits a local physical solution on the
degenerate interval $[0,0]$:

\begin{theorem}[\lean{exists_isLocalPhysicalSolution_zero_time}]
For every $\nu$ and every smooth divergence-free $u_0$ of finite energy, the affine path
$u(t)=u_0+t\,v_0$ with $v_0=\nu\Delta u_0-(u_0\cdot\nabla)u_0$ and zero pressure satisfies
\lean{IsLocalPhysicalSolution ν 0 u₀ u (fun _ => 0)}.
\end{theorem}

The 0.1.5 definitions add the clause \lean{BddAbove (Set.range fun x => ‖vorticity (u t) x‖)} at every
$t\in[0,T]$. With it, data with unbounded initial vorticity never satisfy the bound
(\lean{not_localTetherBound_of_unbounded}, \lean{not_localTetherRate_of_unbounded}); accordingly the closing bound
\lean{LocalTetherClosingBound} and the theorems that consume a bound for all data now range over data with
bounded initial vorticity. Whether the 0.1.4 forms were satisfiable by the junk value alone at some datum (which
would require unbounded vorticity along every local solution at every time) is a question about solutions on
$\R^3$ that is not settled here; the guard removes the possibility.

\subsection{Two-sided differentiability at the endpoint}
The 0.1.4 \lean{LocalTetherRate} required \lean{DifferentiableAt ℝ M s} at every $s\in[0,T]$. At $s=T$ this
depends on the path for $t>T$, where a local solution on $[0,T]$ is unconstrained.

\begin{theorem}[\lean{isLocalPhysicalSolution_escapePath}, \lean{vorticity_sup_norm_escapePath_pos}, \lean{not_localTetherRateV014_zero}, \lean{not_forall_localTetherRateV014}]\label{thm:escape}
Let $u(t)=0$ for $t\le0$ and $u(t)=t^2\,(\sin(t^{-3}x_2),\cos(t^{-3}x_2),0)$ for $t>0$ (\lean{escapePath}). Then $u$
is a local physical solution on $[0,0]$ from $u_0=0$ with zero pressure, and $M(t)=t^{-1}$ for $t>0$ while
$M(0)=0$. Hence $M$ is not differentiable at $0$, \lean{LocalTetherRateV014 0 ν} is false for every $\nu$, and
$\forall u_0\,\nu$, \lean{LocalTetherRateV014 u₀ ν} is false.
\end{theorem}

The last statement is the hypothesis \texttt{hDI} of the 0.1.4 \lean{frohmanian_tether_theorem}. \textbf{The
0.1.4 \lean{frohmanian_tether_theorem} was therefore vacuously true.} The 0.1.5 \lean{LocalTetherRate} asks for
continuity of $M$ on $[0,T]$, a right derivative $d\le CM(s)^2-\kappa''M(s)^3$
(\lean{HasDerivWithinAt M d (Set.Ici s) s}) at every $s\in[0,T)$, and the \lean{BddAbove} clause; it reads $u$ on
$[0,T]$ only, and it holds at $u_0=0$ (\lean{localTetherRate_zero}). The comparison step uses the new
\lean{comparison_ode_stability_Icc_right}, proved from Mathlib's
\lean{image_le_of_liminf_slope_right_lt_deriv_boundary} and \lean{HasDerivWithinAt.liminf_right_slope_le}.

\begin{remark}[The global-path rate; not formalized]
\lean{FrohmanianTetherRate} (global solutions, auxiliary lemmas only; unchanged in 0.1.5) also asks for two-sided
differentiability of $M$ at $s=0$, where \lean{IsPhysicalNSSolution} continues $u$ affinely,
$u(t)=u_0+t\,\partial_tu(0^+)$ for $t<0$. If $\|\omega_0\|$ attains its maximum $M_0>0$ at two points $a,b$ with
$\langle\xi,\partial_t\omega(0^+)\rangle$ equal to $\alpha$ at $a$ and $\beta\neq\alpha$ at $b$ ($\xi=\omega_0/M_0$),
then $M(t)\ge M_0+t\max(\alpha,\beta)+o(t)$ for $t\downarrow0$ and $M(t)\ge M_0+t\min(\alpha,\beta)+O(t^2)$ for
$t\uparrow0$, which is incompatible with differentiability at $0$. Thus \lean{FrohmanianTetherRate} fails at any datum
carrying a global physical solution with such a pair of maxima. No nonzero global finite-energy solution is
formalized, so this is recorded as a mathematical remark.
\end{remark}

\subsection{Statement changes (changelog)}\label{sec:diff}
Re-derived by comparing every declaration header of the Lean sources of 0.1.4 (\texttt{main} at \texttt{8b7bd67})
and 0.1.5. Changed statements:
\begin{itemize}[leftmargin=1.2em]
\item \lean{TetheredLyapunov.LocalTetherBound}: the conclusion \lean{vorticity_sup_norm (vorticity (u t)) ≤ Y}
became \lean{BddAbove (Set.range fun x => ‖vorticity (u t) x‖) ∧ vorticity_sup_norm (vorticity (u t)) ≤ Y}.
\item \lean{TetheredLyapunov.LocalFiniteTimeVorticityBound}: the same change, with \lean{C} in place of \lean{Y}.
\item \lean{TetheredLyapunov.LocalTetherRate}: the clauses
\lean{∀ s ∈ Set.Icc 0 T, DifferentiableAt ℝ M s} and \lean{∀ s ∈ Set.Icc 0 T, deriv M s ≤ C M² − κ'' M³}
became \lean{ContinuousOn M (Set.Icc 0 T)}, \lean{∀ s ∈ Set.Ico 0 T, ∃ d, HasDerivWithinAt M d (Set.Ici s) s ∧ d ≤ C M² − κ'' M³}
and \lean{∀ t ∈ Set.Icc 0 T, BddAbove (Set.range fun x => ‖vorticity (u t) x‖)}.
\item \lean{ClayStatements.LocalTetherClosingBound}: new premise \lean{BddAbove (Set.range fun x => ‖vorticity u₀ x‖)}.
\item \lean{ClayStatements.frohmanian_conditional_regularity_of_closing_bound}: new argument
\lean{(hω₀ : BddAbove (Set.range fun x => ‖vorticity u₀ x‖))}.
\item \lean{ConditionalRegularity.frohmanian_tether_theorem}: \lean{hcont : ∀ u₀ ν, BKMGlobalContinuation u₀ ν}
became \lean{∀ u₀ ν, 0 < ν → BKMGlobalContinuation u₀ ν}; \lean{hDI : ∀ u₀ ν, LocalTetherRate u₀ ν} became
\lean{∀ u₀ ν, 0 < ν → BddAbove (Set.range fun x => ‖vorticity u₀ x‖) → LocalTetherRate u₀ ν}; clause (2) has the
new premise \lean{BddAbove (Set.range fun x => ‖vorticity u₀ x‖)}.
\end{itemize}
Unchanged text, changed meaning through the definitions above: \lean{BKMGlobalContinuation} (its premise
\lean{LocalFiniteTimeVorticityBound} is stronger, so the hypothesis is weaker), \lean{conditional_regularity},
\lean{conditional_regularity_of_local_bound}, both \lean{conditional_regularity_of_open_hyps},
\lean{localFiniteTimeVorticityBound_of_localTetherBound}, \lean{localTetherBound_of_localTetherRate},
\lean{localTetherBound_zero}, \lean{bkmGlobalContinuation_zero}, \lean{conditional_regularity_nonvacuous_at_zero},
\lean{localTetherBound_vacuous_of_periodic}, \lean{bkmGlobalContinuation_vacuous_of_periodic}. Added: the modules
\lean{TetherImpliesBKM} and \lean{TetherDefinitionAudit}, \lean{localTetherRate_zero} and
\lean{comparison_ode_stability_Icc_right}. No other declaration header changed and none was removed. The periodic
statements of Theorem~\ref{thm:mainA} are unchanged from 0.1.4.

